\MWChapter{Semantic architecture and certificates}{语义架构与证书}
\label{ch:architecture}

\section{Layers of the theory / 理论层次}

\begin{table}[H]
\centering
\caption{The semantic layers and their roles in the current research branch.}
\label{tab:components}
\begin{tabularx}{\textwidth}{@{}L{0.23\textwidth}L{0.31\textwidth}Y@{}}
\toprule
\rowcolor{MWSoft!70}
\textbf{Layer} & \textbf{Responsibility} & \textbf{What it does not establish by itself} \\
\midrule
Free SMC presentation & Legal typed sequential, parallel, and symmetric composition & A concrete scheduler or a projection's operational semantics \\
Typed open hypergraph / DPOI & Add, replace, reconnect, or delete a collaboration subgraph under stated boundary and gluing conditions & Global confluence, every proposed match, or every product rule \\
Projection certificates & Connect an identified source event to target transitions and replay evidence & A universal product instance without product-supplied certificates \\
Late-\(\pi\) layer & Sessions, request/accept, mobility, alpha conversion, and structural laws & A released channel API or unrestricted runtime protocol claim \\
Petri layer & Ordered interfaces and individual-token provenance & Boundedness, liveness, or numerical caps without net-level checks \\
Feedback / FMS boundary & Evidence closure, terminal classes, and the selected D1-A semantic route & A single unified source-paper FMS model or unrestricted full abstraction \\
\bottomrule
\end{tabularx}
\end{table}

\section{Identified reconfiguration / 具备身份的重构}

For a configuration \(G\), a concrete rewrite event is represented schematically as
\begin{equation}
  e=(\rho,m,\delta),\qquad G\xrightarrow{e}G',
  \label{eq:identified-event}
\end{equation}
where \(\rho\) is the rule, \(m\) the chosen match, and \(\delta\) retains derivation,
representative, freshness, or allocation information. This is the minimum information needed to
distinguish two applications of the same named rule and to define replay at the selected equality
or isomorphism \citep{cantiluneSpec}.

\section{Projection consistency target / 投影一致性目标}

For each useful reading \(P_i\), the intended condition is
\begin{equation}
  G\xrightarrow{e}G'\quad\Longrightarrow\quad
  P_i(G)\xrightarrow{\Phi_i(e)}P_i(G'),
  \qquad i\in\{\mathrm{DAG},\mathrm{Petri},\pi,\mathrm{morphism}\}.
  \label{eq:projection-consistency}
\end{equation}
The implication is a proof obligation, not a visualization convention. It must include the target
event relation and appropriate preservation, reflection, replay, and terminal-classification
clauses. The morphism reading is the identity view; the remaining readings require explicit
certificates \citep{cantiluneRFCOne,cantiluneRFCTwo}.

\begin{MWCodeBox}[Schematic event-certificate record]
\begin{lstlisting}[style=moonweave]
source event:     operation id, rule id, match, derivation, epoch
projection map:   target event(s), correspondence witness, replay identity
terminal view:    success | external wait | deadlock | productive-infinite
provenance:       source commit, manifest entry, build-evidence reference
\end{lstlisting}
\end{MWCodeBox}

\section{Admission boundary / 准入边界}

Dynamic graph change is intentionally constrained. The theory records structural, freshness,
dangling, quiescence, resource, authorization, and replay conditions as applicable certificates.
A future node or model may propose a new subgraph or channel, but the proposal is not equivalent
to an accepted reconfiguration. This is the formal analogue of a visible control boundary, not a
claim about an already deployed policy service.
